% =============================================================================
%  2-Desarrollo.tex — Cláusulas del acuerdo
% =============================================================================

\section{Cláusulas}

% ── Primera ───────────────────────────────────────────────────────────────────
\clausula{Primera}{Objeto}

El presente acuerdo tiene por objeto establecer los compromisos de Las Partes para respetar la identidad de las personas, acompañarse en el aprendizaje sobre cómo dirigirse a ellas con respeto, y fijar límites y un procedimiento para las bromas que no se entiendan o puedan resultar ofensivas.

% ── Segunda ───────────────────────────────────────────────────────────────────
\clausula{Segunda}{Respeto a la identidad de las personas}

\begin{enumerate}[label=\arabic*., leftmargin=*, itemsep=0.3em]
  \item Las Partes se comprometen a respetar la identidad de las personas, en especial de la comunidad LGBTQ+ y sus asociados, y en general de cualquier otra identidad.
  \item Este compromiso aplica en ambos sentidos: de Rosario hacia Galileo y de Galileo hacia Rosario, así como hacia las personas con quienes cualquiera de ellas conviva.
  \item Las Partes se abstendrán de expresiones ofensivas, burlas o descalificaciones sobre la identidad de cualquier persona, esté o no presente.
  \item Respetar una identidad no obliga a compartirla. Las Partes pueden tener opiniones distintas, siempre que las expresen sin invalidar ni ridiculizar a las personas.
\end{enumerate}

% ── Tercera ───────────────────────────────────────────────────────────────────
\clausula{Tercera}{Explicación anticipada y acompañamiento}

\begin{enumerate}[label=\arabic*., leftmargin=*, itemsep=0.3em]
  \item \textbf{Explicación anticipada.} Cuando Las Partes vayan a convivir con una persona o comunidad cuya identidad requiera un trato particular, la Parte que conozca el tema dará, de manera previa y dentro de lo posible, una explicación anticipada sobre:
  \begin{enumerate}[label=\alph*), itemsep=0.2em]
    \item Lo que puede resultar ofensivo para esa persona o comunidad.
    \item Las maneras respetuosas de dirigirse a ella, por ejemplo, los nombres, términos o formas de trato que prefiere.
    \item Los temas o comentarios que conviene evitar.
  \end{enumerate}

  \item \textbf{Disposición a aprender.} La Parte que desconozca el tema se compromete a escuchar la explicación, a preguntar cuando tenga dudas y a no tomarla como un ataque o un reproche.

  \item \textbf{Señalamiento respetuoso.} Si durante una convivencia una de Las Partes realiza una expresión ofensiva por desconocimiento, la otra Parte se lo señalará con respeto, de preferencia en privado, explicando por qué puede resultar ofensiva.

  \item \textbf{Buena fe.} Un error cometido por desconocimiento no se presume mala intención. La Parte señalada se compromete a reconocerlo y a corregirlo en adelante, sin justificar la expresión ni minimizar el sentir de la otra Parte.
\end{enumerate}

% ── Cuarta ────────────────────────────────────────────────────────────────────
\clausula{Cuarta}{Bromas sobre la muerte o el suicidio}

\begin{enumerate}[label=\arabic*., leftmargin=*, itemsep=0.3em]
  \item \textbf{Prohibición.} Ninguna de Las Partes incluirá a la otra en bromas sobre la muerte o el suicidio. Esta obligación recae de manera predominante en Galileo hacia Rosario, pero no exime a Rosario de la misma obligación hacia Galileo.

  \item \textbf{Bromas sobre la propia persona.} Cada Parte puede hacer bromas sobre la muerte o el suicidio siempre y cuando traten únicamente de su propia persona y no de la otra Parte.

  \item \textbf{Pregunta en serio.} Si una broma sobre la propia muerte o el suicidio genera preocupación genuina, la otra Parte puede preguntar con seriedad si detrás de ella existe algo real. La Parte que hizo la broma debe responder con honestidad, sin evadir la pregunta con otra broma.

  \item \textbf{Riesgo a la integridad.} Si de la respuesta se desprende un riesgo para la salud o la integridad de cualquiera de Las Partes, se aplica la Cláusula Quinta del Acuerdo 1: la prioridad es atender a quien lo necesite y, de ser necesario, pedir ayuda a terceros o a servicios de emergencia.
\end{enumerate}

% ── Quinta ────────────────────────────────────────────────────────────────────
\clausula{Quinta}{Procedimiento ante una broma no entendida u ofensiva}

Cuando una broma no se entienda o pueda resultar ofensiva para cualquiera de Las Partes, se seguirán los siguientes pasos:

\begin{enumerate}[label=\textbf{Paso \arabic*.}, leftmargin=*, itemsep=0.3em]
  \item \textbf{Pausa de aclaración.} La Parte que no entienda la broma o se sienta incómoda puede detenerla para preguntar \fraseclave{¿Es broma?}. La Parte que hizo la broma debe responder con respeto, sin hacer más grande el chiste ni la situación.

  \item \textbf{Evaluación del límite.} La Parte que preguntó decide si la broma entra dentro de su límite personal. Si es así, puede seguir con el chiste. Si no lo es, se pasa al Paso 3.

  \item \textbf{Rechazo respetuoso.} Si la broma pasa el límite personal, la Parte que preguntó puede rechazarla. Para ello:
  \begin{enumerate}[label=\alph*), itemsep=0.2em]
    \item Termina la broma de manera respetuosa, sin ninguna acción violenta, ya sea física, psicológica o social.
    \item Informa a la otra Parte que la broma pasó su límite y le aclara su sentir y su pensamiento.
    \item Las Partes cambian de tema sin mayor problema.
  \end{enumerate}
\end{enumerate}

La Parte que hizo la broma no debe insistir en ella, repetirla ni reprochar el rechazo. Rechazar una broma no constituye un ataque, y aceptarla en una ocasión no obliga a aceptarla en otra.

\begin{table}[H]
  \centering
  \small
  \begin{tabularx}{\textwidth}{@{} l >{\raggedright\arraybackslash}p{3.6cm} X @{}}
    \toprule
    \textbf{Paso}            & \textbf{Quién actúa}          & \textbf{Regla} \\
    \midrule
    1. Pausa de aclaración   & Quien no entiende la broma    & Pregunta \enquote{¿Es broma?}; la otra Parte responde con respeto y sin agrandar el chiste. \\
    2. Evaluación del límite & Quien preguntó                & Decide si la broma entra en su límite: si entra, se sigue con el chiste. \\
    3. Rechazo respetuoso    & Quien preguntó                & Termina la broma sin violencia, aclara su sentir y se cambia de tema. \\
    \bottomrule
  \end{tabularx}
  \caption{Resumen del procedimiento ante una broma no entendida u ofensiva.}
  \label{tab:bromas}
\end{table}

% ── Sexta ─────────────────────────────────────────────────────────────────────
\clausula{Sexta}{Relación con el Acuerdo 1}

\begin{enumerate}[label=\arabic*., leftmargin=*, itemsep=0.3em]
  \item Si una broma o una expresión ofensiva escala a una discusión que desborde a cualquiera de Las Partes, esta puede solicitar una pausa activa conforme a la Cláusula Cuarta del Acuerdo 1.
  \item Las excepciones por integridad, seguridad o salud de la Cláusula Quinta del Acuerdo 1 aplican también a este acuerdo.
\end{enumerate}

% ── Séptima ───────────────────────────────────────────────────────────────────
\clausula{Séptima}{Vigencia y modificaciones}

\begin{enumerate}[label=\arabic*., leftmargin=*, itemsep=0.3em]
  \item El presente acuerdo entra en vigor en la fecha de su firma y tiene vigencia indefinida.
  \item Cualquiera de Las Partes puede solicitar su revisión. Toda modificación debe hacerse de mutuo acuerdo y por escrito, mediante una adenda o un nuevo acuerdo.
\end{enumerate}
